\section*{Current Progress}
\subsection*{Literature Review, Technical Familiarization and other Preliminaries}

\subsection*{Implementation of Dirty Access Tracking in UVM Driver}

\paragraph{Workflow.}
A tracking epoch begins when a privileged user-space tool writes to
\texttt{dirty\_pids\_start\_track}.
This initialises a fresh per-process table and immediately invalidates all
current GPU page-table entries by unmapping them, so that every subsequent GPU
access generates a fresh page fault.
Each write fault is intercepted in the driver's fault-servicing path, which
records the faulting page in the table along with a nanosecond-resolution
timestamp before re-establishing the mapping and replaying the faulting
instruction.
At any point during the epoch the tool can read \texttt{dirty\_pages} to obtain
the accumulated set of written pages; reads are non-destructive, so the table
is not consumed by the query.
The epoch ends when the tool writes to \texttt{dirty\_pids\_stop\_track},
which frees all recorded state. Re-starting tracking destroys any existing
table before creating a fresh one, cleanly beginning a new epoch.

\paragraph{Data Structures.}
The tracking state is organised as a two-level structure defined in
\texttt{uvm\_common.h}.

\begin{enumerate}
    \item \texttt{struct dirty\_page\_info} - a per-page record holding the
          page number, the \texttt{ktime\_get\_ns()} timestamp of the first
          observed write fault, and the PID of the faulting process.
    \item \texttt{struct uvm\_dirty\_page\_table} - a per-process container
          holding an \texttt{xarray} that maps page numbers to
          \texttt{dirty\_page\_info} pointers.
    \item A global \texttt{xarray} (\texttt{pid\_to\_page\_table}) maps each
          tracked process's \texttt{tgid} to its \texttt{uvm\_dirty\_page\_table}.
\end{enumerate}

This two-level design provides natural per-process isolation. Entries are recorded under a
\emph{first-write-wins} policy: an existing entry is never overwritten, so the
stored timestamp always reflects the earliest observed write to that page within
the current epoch. The faults may be serviced by different CPU-side driver threads, 
thus the timestamp may not be strictly monotonic, depending upon the driver-side service thread
recording the same.

\paragraph{APIs Exposed.}
User-space interacts with the tracker exclusively through five
\texttt{procfs} entries created under \texttt{/proc/driver/nvidia-uvm/}:

\begin{enumerate}
    \item \texttt{dirty\_pids\_start\_track} (write) - initialises (or
          re-initialises) the tracking table for \texttt{current->tgid} and
          triggers GPU-mapping invalidation to start a new epoch.
    \item \texttt{dirty\_pids\_stop\_track} (write) - destroys the table for
          \texttt{current->tgid}, freeing all recorded entries.
    \item \texttt{dirty\_pid\_to\_query} (write) - sets the PID whose table
          is consulted on subsequent \texttt{dirty\_pages} reads, decoupling
          the querying process from the tracked process.
    \item \texttt{dirty\_pages} (read) - iterates the selected process's
          table and emits one line per recorded page in the format
          \texttt{0x\textit{addr}\ \textit{timestamp\_ns}\ \textit{pid}}.
          Reads are \emph{non-destructive}: entries remain in the table and
          can be queried repeatedly within the same epoch.
    \item \texttt{dirty\_range} (write) - accepts a pair of hexadecimal
          addresses \texttt{0x\textit{start}\ 0x\textit{end}} (end exclusive)
          to restrict subsequent \texttt{dirty\_pages} reads to a specific
          virtual address window. The range is protected by a spinlock to
          allow concurrent updates safely. An inverted or zero-length range
          produces an empty result rather than undefined behaviour.
\end{enumerate}

\paragraph{Internal Functions and Implementation Sites.}
The implementation is spread across several driver files, each responsible for
a distinct concern:

\begin{enumerate}
    \item \texttt{uvm\_common.c} - the primary implementation file.
          Contains \texttt{uvm\_dirty\_page\_table\_init} /
          \texttt{uvm\_dirty\_page\_table\_destroy} (table lifecycle),
          \texttt{uvm\_dirty\_page\_table\_record} (fault recording, using
          \texttt{GFP\_ATOMIC} as it executes in a non-sleepable context),
          \texttt{uvm\_dirty\_page\_table\_lookup}, and the full
          \texttt{procfs} handler implementations including
          \texttt{uvm\_dirty\_procfs\_init}.
    \item \texttt{uvm\_gpu\_replayable\_faults.c} - the write-fault hook.
          After standard fault servicing, checks
          \texttt{uvm\_dirty\_tracking\_active\_for\_pid(va\_block->creator\_pid)}
          and calls \texttt{uvm\_dirty\_page\_table\_record} for
          \texttt{UVM\_FAULT\_ACCESS\_TYPE\_WRITE} faults.
    \item \texttt{uvm\_va\_space.c} - implements
          \texttt{uvm\_dirty\_invalidate\_all\_gpu\_mappings}, which walks all
          live VA spaces and blocks to unmap GPU PTEs at epoch start. Exposes
          \texttt{uvm\_va\_space\_dirty\_init} to register this function via
          the \texttt{uvm\_dirty\_invalidate\_fn} function pointer at module
          load.
    \item \texttt{uvm.c} - top-level integration; calls
          \texttt{uvm\_dirty\_procfs\_init} and
          \texttt{uvm\_va\_space\_dirty\_init} during module initialisation.
\end{enumerate}

\subsection*{Initial Performance Evaluation}
